\section{A frame character beyond the metabelian action}
\label{sec:frame-witness}

The quadratic beta characters are detected by metabelian homology.
The following construction gives a different kind of element: its
normalized conjugator lies in the third derived subgroup, while one
of its torsion frames has a nontrivial central value modulo twelve.
The construction uses the relative form of
Theorem~\ref{thm:peripheral-lifting}, with the closed normal subgroup
$U=F''$; no openness assumption on $U$ is needed.

Keep the modular marking $\Gamma=\langle s,r\mid s^2=r^3=1\rangle$,
$\tau=sr$, $x=\tau^{-2}$, and $y=sxs$. Let
\[
 q:\Gamma\longrightarrow
 Q_{12}=\operatorname{SL}_2(\mathbb Z/12\mathbb Z)/\{\pm I\}
\]
be the congruence quotient with the signed pure images
\[
 X=\begin{pmatrix}1&2\\0&1\end{pmatrix},\qquad
 Y=\begin{pmatrix}1&0\\-2&1\end{pmatrix}.
\]
These arise from $s=\left(\begin{smallmatrix}0&1\\-1&0\end{smallmatrix}\right)$
and $\tau=\left(\begin{smallmatrix}1&-1\\0&1\end{smallmatrix}\right)$,
with matrices taken modulo scalar sign. Put $\zeta=[5I]\in Q_{12}$.

\begin{theorem}\label{thm:deep-frame-witness}
There is a multiplier-one gentle element $e=(1,f)$ such that
\[
 f\in F''',\qquad \rho_2(e)=0.
\]
Its modular automorphism $\Phi$ fixes $\tau$ and induces the identity
under $q$. Its unique pure torsion frames, defined by
\[
 \Phi(s)=\beta s\beta^{-1},\qquad
 \Phi(r)=ArA^{-1},
\]
satisfy
\[
 q(\beta)=1,\qquad q(A)=\zeta.
\]
In particular $e$ acts identically on $F/F'''$, and every character
factoring through its metabelian action vanishes on $e$.
\end{theorem}

\begin{proof}
We first calculate the finite image needed for the lift. Set
$z=(xy)^{-1}$ and
\[
 a=[x^2,y^2],\qquad b=[y^2,z^2],\qquad w=[a,b]\in F''.
\]
Direct multiplication gives
\[
 \begin{array}{c|c|c|c}
  &q(a)&q(b)&q(w)\\ \hline
  \bmod3&
    \begin{pmatrix}1&1\\1&2\end{pmatrix}&
    \begin{pmatrix}0&2\\1&0\end{pmatrix}&-I\\[3pt]
  \bmod4&I&I&I\\[3pt]
  \bmod12&
    \begin{pmatrix}1&4\\4&5\end{pmatrix}&
    \begin{pmatrix}9&8\\4&9\end{pmatrix}&5I.
 \end{array}
\]
Modulo three, $X,Y$ generate $\operatorname{SL}_2(\mathbb F_3)$.
The two matrices in the first row generate its normal quaternion
subgroup of order eight: each squares to $-I$, they anticommute,
and conjugation by $X,Y$ preserves this subgroup. The quotient has
order three, so the derived subgroup is precisely this quaternion
group; its derived subgroup is $\{I,-I\}$. Modulo four, $X,Y$
commute and have order two. These calculations give
\[
 q(F'')\subset\{1,\zeta\},
\]
and the word $w$ gives equality. The element $\zeta$ is central
and has order two.

Put $U=F''$. For a prime $p$, let $O^p(U)$ be the intersection of
the kernels of all continuous homomorphisms from $U$ to finite
$p$-groups, and set $\Gamma_p=\Gamma/O^p(U)$.
The subgroup $O^p(U)$ is characteristic in $U$ and hence normal in
$\Gamma$. Write
\[
 R=\bigcap_{n\geq1}\gamma_n(U).
\]
There is a natural isomorphism
\begin{equation}
 \Gamma/R\simeq
   \prod_{p,\,\Gamma/U}\Gamma_p,
 \label{eq:relative-pronilpotent-product}
\end{equation}
where the right side is the fibre product with a common image in
$\Gamma/U$. Here are the details for an arbitrary closed $U$.
Its maximal pronilpotent quotient is the product of its maximal
pro-$p$ quotients: every finite subdirect product of groups of
pairwise coprime orders is the full product, and compactness passes
this surjectivity to the inverse limit. The kernel is the
intersection of the lower central series. Indeed a finite
nilpotent quotient kills some lower central term; conversely an
open subgroup containing the intersection contains some term by
compactness of a descending closed family. To obtain
\eqref{eq:relative-pronilpotent-product}, lift the common
$\Gamma/U$ coordinate to $\Gamma$ and then use the just-proved
surjectivity on $U$. The natural map is a continuous bijection
between compact Hausdorff groups and is therefore an isomorphism.

Since $q(U)=\langle\zeta\rangle$ is a two-group, the map $q$
factors through $\Gamma_2$. Let
\[
 J=\operatorname{im}(\Gamma\longrightarrow Q_{12}\times\Gamma/U).
\]
The map $\Gamma_2\to J$ has pro-two kernel. The element
$(\zeta,1)$ belongs to $J$, by the word $w$, and centralizes the
image of $r$. Choose a preimage $c_0$ in the image $U^{(2)}$ of
$U$ in $\Gamma_2$. The elements $c_0rc_0^{-1}$ and $r$ have
order three and map to the same generator of order three in $J$.
They are complements to a pro-two kernel in the preimage of that
cyclic group. Coprime complement conjugacy in every finite quotient,
followed by compactness, gives a correction in the kernel that
makes them equal. Thus there is
\[
 c\in C_{\Gamma_2}(r)\cap U^{(2)},\qquad q(c)=\zeta.
\]
For clarity, the conjugacy is by a kernel element, so it retains
the image $(\zeta,1)$; equality of the complements implies equality
of the chosen generators because both map to the same element of
order three.

In the product \eqref{eq:relative-pronilpotent-product}, let
$\overline b_r$ have coordinate $c$ at two and identity at every
other prime. It centralizes the image of $r$. Put
$\overline b_s=\overline b_{\tau^{-1}}=1$.
The identity automorphism of $\Gamma/R$ is consequently a positive
peripheral automorphism with these specified conjugators.
Theorem~\ref{thm:peripheral-lifting} supplies a full peripheral
automorphism $\Phi$ with all three specified conjugator images. Its cusp
conjugator lies in $R$; the allowed inner correction by that
conjugator makes $\Phi(\tau)=\tau$ and leaves all their images
in $\Gamma/R$, and hence in $Q_{12}$, unchanged.

The resulting conjugator of $s$ belongs to $R$, and that of $r$
belongs to $U$. Both are pure. The torsion-centralizer theorem for
a free profinite product gives $C_\Gamma(s)=\langle s\rangle$
and $C_\Gamma(r)=\langle r\rangle$, so these pure frames are
unique. They are the $\beta,A$ in the statement, with
$q(\beta)=1$ and $q(A)=\zeta$.

Since $\Phi$ is the identity modulo $R$ and fixes $\tau$, write
\[
 \Phi(s)=sf,\qquad \Phi(r)=f^{-1}r,
       \qquad f=s\Phi(s)\in R.
\]
The residual $R$ is contained in $U'$ because abelian quotients are
nilpotent. Hence $f\in U'=F'''$. The torsion relations give the
two cycle identities, and restriction to $F$ is
$x\mapsto x$, $y\mapsto f^{-1}yf$. There is no additional
abelianization correction because $f$ is already in $F'$.
This is the desired multiplier-one gentle element.

The frame definition of $\rho_2$ reads the abelianization of
$\beta^{-1}$. As $\beta\in R\subset F'''$, this abelianization
is zero, giving $\rho_2(e)=0$. The map $q$ kills $R$, so
$q\Phi=q$. Finally $f\in F'''$ gives literal identity on
$F/F'''$, which includes its metabelian quotient.
\end{proof}

\begin{corollary}\label{cor:frame-character-nonzero}
The character $\Delta$ of Theorem~\ref{thm:global-frame-character}
is nonzero on $N_0$ and is outside the span of all continuous quadratic
characters factoring through the metabelian action. In particular
$\Delta$ and the characters $\vartheta_p$ of
Theorem~\ref{thm:quadratic-beta-family} are linearly independent.
\end{corollary}

\begin{proof}
For the element of Theorem~\ref{thm:deep-frame-witness}, the two
signed matrices are $q_{12}(A)=5I$ and $q_{12}(\beta)=I$, because
the pure image $E_{12}$ embeds in the projective quotient. Hence the
frame matrices modulo three are $q_3(A)=-I$ and
$q_3(\beta)=I$. Definition~\eqref{eq:frame-sign-definition} gives
$a_3(e)=1$, $a_2(e)=0$. Thus $\xi(e)=5\pmod{12}$ and
$\Delta(e)=\delta(e)=1$. Its metabelian action is the identity,
so every character factoring through that action vanishes on it.
Evaluating a finite linear relation on $e$ first kills the coefficient
of $\Delta$; independence of the $\vartheta_p$ kills all the others.
\end{proof}
